\uuid{tQGU}
\titre{Équations complexes}
\auteur{bodin ; hanani ; mzari}
\organisation{exo7}
\licence{CC BY-NC-SA 4.0}
\datecreate{2020-01-20}
\niveau{L1}
\module{Algèbre}
\chapitre{Nombres complexes}
\sousChapitre{Racine carrée, équation du second degré}
\difficulte{3}
\theme{Équations, équations complexes}
\exo7qcmid{lille-1.nombres-complexes.027}

\contenu{
\begin{question}
\qclassification{104.02,104.03,104.04}
\qtags{motcle=équations complexes}
Soit $n$ un entier $\ge 2$, $z_1,z_2, \dots, z_n$ les racines $n$-ièmes de l'unité. Quelles sont les assertions vraies ?
\begin{answers}
    \good{$z^n-1=(z-z_1)(z-z_2)\dots (z-z_n)$}

    \good{$z_1.z_2, \dots z_n = (-1)^{n-1}$ }

    \bad{$z_1+z_2+ \dots + z_n = 1$ }

    \good{$z_1+z_2+ \dots + z_n = 0$ }

\end{answers}
\begin{explanations}
$z_1,z_2, \dots, z_n$ sont les racines dans $\Cc$ du polynôme $P(X) =X^n-1$, donc $P(X)=(X-z_1)(X-z_2)\dots (X-z_n)$.
On examine le coefficient de $X^{n-1}$ et le coefficient constant.
\end{explanations}
\end{question}
}
